\chapter{PRIME Ledger --- Utility and Structural Tags}
\label{chap:PRIME_Ledger}

\section{Purpose}

PRIME tags serve operational and structural utility functions across the
Tagging System and GEL. They do not represent conceptual meaning, but rather
organizational and functional metadata.

\section{Scope}

These tags include:
\begin{itemize}
    \item anchors,
    \item ribbons and markers,
    \item structural hints,
    \item indexing utilities,
    \item workflow or processing flags.
\end{itemize}

\section{Schema}

\[
T_{\text{PRIME}} = \langle key,\, gloss,\, bindings,\, constraints \rangle
\]

Examples:
\begin{itemize}
    \item \texttt{prime:anchor}
    \item \texttt{prime:important}
    \item \texttt{prime:index}
    \item \texttt{prime:role}
    \item \texttt{prime:workflow}
\end{itemize}

\section{Class Binding Notes}

\begin{itemize}
    \item PRIME tags rarely use C/P unless linking to identity invariants.
    \item S-class allowed for structural indexing.
    \item O-class allowed for narrative workflow hints.
\end{itemize}

\section{Constraints}

\begin{itemize}
    \item PRIME tags MUST NOT carry semantic content.
    \item They must not override W5H meaning.
    \item PRIME tags must never be used as substitutes for conceptual tags 
          (SOCI/PSY/ONTO/etc.).
\end{itemize}
